% 2009 年辽宁卷理科第 17 题：海面 A、C 两处测量两座灯塔塔顶 B、D 的仰角。
% 来源 PDF：普通高考/2009/2009辽宁理.pdf，SHA256 5ae9708acfe7def5bddd60a34d1474cfa2e2c20a39d130255e0e4393093accad
% 原页：第 1 页（A3 三栏）第 3 栏中下部，印刷页 537。
% 重绘依据：水面直线过 A、C；A 处测得 B 的仰角 75°（向左上方）、D 的仰角 30°（向右上方）；
%   C 处测得 B、D 的仰角均为 60°（分别向左上方、右上方）。原图岛屿与灯塔为写意画法，
%   此处按同一几何关系改画为示意灯塔，仰角标注位置与原页一致。
\begin{tikzpicture}[scale=0.86, semithick, >=Stealth,
  every node/.style={inner sep=1.5pt, font=\small}]
  % 水面
  \draw (-4.9,0) -- (7.4,0);
  % 两座灯塔（岛与塔身）
  \draw (-4.35,0) -- (-3.75,0.75) -- (-3.35,1.35) -- (-3.0,1.75) -- (-2.55,0.8) -- (-2.1,0);
  \draw (-3.0,1.75) -- (-3.0,6.6);
  \draw (-3.12,6.45) rectangle (-2.88,6.78);
  \draw (5.25,0) -- (5.6,0.55) -- (5.9,1.0) -- (6.25,0.45) -- (6.6,0);
  \draw (5.9,1.0) -- (5.9,2.4);
  \draw (5.78,2.26) rectangle (6.02,2.56);
  % 视线
  \draw (0,0) -- (-3.0,6.6);
  \draw (4,0) -- (-3.0,6.6);
  \draw (0,0) -- (5.9,2.4);
  \draw (4,0) -- (5.9,2.4);
  % 仰角弧与标注
  \draw (0.72,0) arc (0:22.1:0.72);
  \node at (1.30,0.27) {$30^\circ$};
  \draw (-0.2973,0.6558) arc (114.4:180:0.72);
  \node at (-0.86,0.54) {$75^\circ$};
  \draw (4.72,0) arc (0:51.6:0.72);
  \node at (4.86,0.30) {$60^\circ$};
  \draw (3.4760,0.4936) arc (136.7:180:0.72);
  \node at (2.88,0.58) {$60^\circ$};
  % 顶点标注
  \node[above] at (-3.0,6.86) {$B$};
  \node[right] at (6.06,2.48) {$D$};
  \node[below] at (0,-0.10) {$A$};
  \node[below] at (3.985,-0.10) {$C$};
\end{tikzpicture}
